\documentclass[11pt]{article}
\usepackage[T1]{fontenc}
\usepackage[utf8]{inputenc}
\usepackage{lmodern}
\usepackage[margin=1in]{geometry}
\usepackage{amsmath,amssymb,amsthm,mathtools,booktabs,array}
\usepackage{microtype}
\usepackage{xcolor}
\usepackage{hyperref}
\usepackage{fancyhdr}
\pagestyle{fancy}
\fancyhf{}
\fancyhead[L]{\small Alec Kriebel}
\fancyhead[R]{\small Kourovka Notebook 16.63}
\fancyfoot[C]{\thepage}
\setlength{\headheight}{14pt}
\hypersetup{colorlinks=true,linkcolor=blue!45!black,citecolor=blue!45!black,urlcolor=blue!45!black,
 pdftitle={An explicit odd-order p-group with as many automorphisms as elements},
 pdfauthor={Alec Kriebel},
 pdfsubject={Computer-assisted solution to Kourovka Notebook problem 16.63}}
\newtheorem{theorem}{Theorem}[section]
\newtheorem{lemma}[theorem]{Lemma}
\newtheorem{proposition}[theorem]{Proposition}
\newtheorem{corollary}[theorem]{Corollary}
\theoremstyle{definition}\newtheorem{definition}[theorem]{Definition}
\theoremstyle{remark}\newtheorem{remark}[theorem]{Remark}
\DeclareMathOperator{\Aut}{Aut}
\DeclareMathOperator{\Der}{Der}
\DeclareMathOperator{\End}{End}
\DeclareMathOperator{\Hom}{Hom}
\DeclareMathOperator{\ad}{ad}
\DeclareMathOperator{\BCH}{BCH}
\DeclareMathOperator{\GL}{GL}
\DeclareMathOperator{\PGL}{PGL}
\DeclareMathOperator{\rank}{rank}
\newcommand{\Zp}{\mathbb Z_p}
\newcommand{\Qp}{\mathbb Q_p}
\newcommand{\Fp}{\mathbb F_p}
\newcommand{\Span}{\operatorname{span}}
\newcommand{\EE}{\mathcal E}
\newcommand{\TT}{\mathcal T}
\newcommand{\FF}{\mathcal F}
\setlength{\emergencystretch}{2em}
\AtBeginDocument{\setlength{\abovedisplayshortskip}{7pt}}
\title{An explicit odd-order $p$-group with as many\\automorphisms as elements\\[5pt]
\large A computer-assisted solution to Kourovka Notebook problem 16.63}
\author{Alec Kriebel\\[3pt]
\small Independent researcher\\
\small\href{https://orcid.org/0009-0001-9320-500X}{ORCID: 0009-0001-9320-500X}\\
\small\href{mailto:me@aleckriebel.com}{me@aleckriebel.com}}
\date{Version 1.1.0 --- 15 September 2026}
\begin{document}
\maketitle
\begin{abstract}
We construct a finite group $G$ for which
\[
 p=1009,\qquad |G|=|\Aut(G)|=1009^{52359}.
\]
The group is the Baker--Campbell--Hausdorff group of an explicitly specified
rank-$31$ Lie ring modulo $1009^{1689}$. Its nilpotency class is at most $846<1009$.
An asymmetric lattice forces every automorphism of the finite Lie ring into an
integral logarithm--exponential neighborhood. This gives a bijection between the
\emph{full} automorphism set and the kernel of a linear derivation matrix.
Exact independently implemented Python and C++ computations establish the
Smith valuations of that matrix. The resulting computer-assisted proof answers the existence question
affirmatively. The construction and all exact certificates are supplied
with separately implemented verification software.
\end{abstract}
\medskip
\noindent\textbf{Mathematics Subject Classification (2020):} 20D45, 20D15, 17B40.\\
\textbf{Keywords:} finite $p$-groups; automorphism groups; Lie lattices;
Lazard correspondence; Smith normal form.

\medskip
\noindent\textbf{AI disclosure.} The construction, proof development, and initial
verification software were generated using ChatGPT-6 Astra Pro. Separate Codex
sessions performed additional AI-assisted reviews and exact computational checks.
These checks do not constitute external human peer review or proof-assistant
formalization.

\medskip
\noindent\textbf{Archive DOI for this edition:}
\href{https://doi.org/10.5281/zenodo.22770864}{10.5281/zenodo.22770864}.
Source and certificate availability is detailed in Section~\ref{sec:reproducibility}.

\clearpage
{\small\tableofcontents}
\clearpage

\section{Statement, sources, and scope}
Problem 16.63, attributed to D. MacHale in the September 2026 version of the
\emph{Kourovka Notebook}, asks for a nontrivial finite $p$-group of
odd order with $|\Aut(G)|=|G|$ \cite{KN}. It appears on printed page 102 and
cross-refers to Archive 12.77 on printed page 255. The latter is the separate,
now-disproved divisibility conjecture. The notebook website identifies its
1 September 2026 file as the current update. No published resolution of the
exact equality question was located in the literature searches recorded in the
accompanying ledger.

Gonz\'alez-S\'anchez and Jaikin-Zapirain construct, for every prime, families with
$|\Aut(G)|=O(|G|^{40/41})$ \cite{GSJ}. This explains the negative answer to the
archival divisibility problem; it does not itself give the equality required
here. Shabani-Attar's paper treats particular classes: its Proposition 2.5,
Theorem 2.6 and Theorem 2.8 exclude odd-prime examples of maximal class, with
cyclic Frattini subgroup, and of class two with cyclic centre, respectively
\cite{SA}. None is a general nonexistence theorem.

The Lie-algebra template used below is the rank-$31$ example in
Omirov--Ruan, Theorem 4.8, with six radical summands \cite{OR}. Their arXiv
record has version 1, dated 6 May 2026. We specify an integral normalization
and establish the required finite-characteristic facts directly; we do not
assume that an unspecified characteristic-zero automorphism theorem reduces
unchanged modulo $p$. The asymmetric lattice and finite full-automorphism
count are developed in this report.

As elementary controls, a nontrivial abelian group of odd order has the
inversion automorphism of order two, so cannot answer the question. For
\[
 H_p=\langle a,b\mid a^{p^2}=b^p=1,\;bab^{-1}=a^{1+p}\rangle,
\]
the normal forms $a^u b^v$ give $|H_p|=p^3$. All automorphisms have images
$a\mapsto a^u b^v$, $b\mapsto a^{pw}b$, where $u$ is a unit modulo $p^2$
and $v,w\in\Fp$. Thus $|\Aut(H_p)|=(p-1)p^3$, not $p^3$.
The exhaustive generating-pair controls give $(27,54)$, $(125,500)$ and
$(343,2058)$ for $(|H_p|,|\Aut(H_p)|)$ at $p=3,5,7$.
This also checks why counting only a Sylow subgroup would give a false positive.

\begin{theorem}\label{thm:main}
There exists a finite group $G$ of order $1009^{52359}$ whose full automorphism
group has order $1009^{52359}$. A portable description is given in
Definitions~\ref{def:B} and \ref{def:A} and equation~\eqref{eq:G}.
\end{theorem}
The prime, order and number of generators are not claimed to be minimal.
The theorem provides one example, not an assertion for every odd prime.

\section{The integral Lie algebra}
Fix $p=1009$. This number is prime, as is also checked by trial division in
the independent verifier. All lattices in the proof are over $\Zp$.
For $m\ge0$, write $V_m$ for the free module of homogeneous binary forms of
degree $m$, with basis $v_i=X^{m-i}Y^i$, $0\le i\le m$.
The standard $\mathfrak{sl}_2$ action is
\begin{equation}\label{eq:action}
 [h,v_i]=(m-2i)v_i,\quad [e,v_i]=i v_{i-1},\quad
 [f,v_i]=(m-i)v_{i+1},
\end{equation}
with out-of-range terms zero, and
$[h,e]=2e$, $[h,f]=-2f$, $[e,f]=h$.

For binary forms $F\in V_m$, $H\in V_n$, define the \emph{raw transvectant}
\begin{equation}\label{eq:transvectant}
 P_r(F,H)=\sum_{k=0}^r(-1)^k\binom rk
 \frac{\partial^r F}{\partial X^{r-k}\partial Y^k}
 \frac{\partial^r H}{\partial X^k\partial Y^{r-k}}.
\end{equation}
No factorial normalization is suppressed. On basis vectors the coefficient
of the vector indexed by $i+j-r$ in $V_{m+n-2r}$ is
\begin{equation}\label{eq:falling}
 \sum_{k=0}^r(-1)^k\binom rk
 (m-i)_{r-k}(i)_k(n-j)_k(j)_{r-k},
\end{equation}
where $(a)_b=a(a-1)\cdots(a-b+1)$, $(a)_0=1$, and $(a)_b=0$ when $b>a\ge0$.

\begin{definition}\label{def:B}
Let $B$ have ordered basis
\[
 h,e,f,\quad U_0,\ldots,U_6,\quad V_0,\ldots,V_4,\quad
 W_0,\ldots,W_6,\quad T_0,T_1,T_2,\quad R_0,\ldots,R_4,\quad Z_0.
\]
Its six radical summands $U,V,W,T,R,Z$ are copies of
$V_6,V_4,V_6,V_2,V_4,V_0$, respectively. In addition to
\eqref{eq:action}, the only nonzero radical brackets, up to skew symmetry, are
\begin{center}
\begin{tabular}{ccc}
\toprule Inputs & Output summand & Map \\ \midrule
$U,U$ & $W$ & $P_3$ \\
$U,V$ & $T$ & $P_4$ \\
$U,W$ & $T$ & $P_5$ \\
$U,R$ & $T$ & $P_4$ \\
$V,V$ & $T$ & $P_3$ \\
$V,R$ & $Z$ & $P_4$ \\
$R,R$ & $T$ & $P_3$ \\
\bottomrule
\end{tabular}
\end{center}
The target of a transvectant is identified with the indicated copy of its
binary-form module. All other brackets are zero.
\end{definition}
This is a complete specification with integer coefficients. In particular,
$\rank B=3+7+5+7+3+5+1=31$.

\begin{lemma}\label{lem:basic}
The displayed bracket is a Lie bracket over $\mathbb Z$. Over $\Zp$, $\Qp$
and $\Fp$ it is perfect and has centre precisely the $Z$ summand. If
$N=U\oplus V\oplus W\oplus T\oplus R\oplus Z$, then
\[
 [N,N]=W\oplus T\oplus Z,\qquad [N,[N,N]]=T,\qquad [N,T]=0.
\]
\end{lemma}
\begin{proof}
The action and transvectants are $\mathfrak{sl}_2$-equivariant. Hence the only
potentially nonzero radical Jacobi expression is the alternating expression
on three $U$ inputs, taking values in $T$. Over the rationals,
\[
 \Lambda^3 V_6\cong V_{12}\oplus V_8\oplus V_6\oplus V_4\oplus V_0,
\]
as follows by multiplying the seven weight characters and taking the third
exterior power. It has no $V_2$ summand. The equivariant Jacobi map to $T=V_2$
is therefore zero. Its coefficients are integers, so this proves the
integer identity. Independently, both implementations check all
$\binom{31}{3}=4495$ basis triples, with integer arithmetic.

The nontrivial simple summands have no $\mathfrak{sl}_2$-fixed vector at
$p=1009$. Thus the centre is contained in $Z$, which is central. Every
nontrivial summand is generated by its brackets with $\mathfrak{sl}_2$, and
$[V,R]=Z$, so $B$ is perfect. The stated lower central terms of $N$ follow
directly from the table. Surjectivity here can be checked on highest weight
vectors; coefficients such as $720$, $8640$ and $576$ are units at $1009$.
\end{proof}

\subsection{Two generators of the ambient lattice}
Set
\begin{equation}\label{eq:xy}
 x=h,\qquad y=e+f+U_1+V_0+R_2,\qquad t=U_0.
\end{equation}
\begin{lemma}\label{lem:generators}
The Lie subalgebra over $\Zp$ generated by $x,y$ is $B$.
\end{lemma}
\begin{proof}
The four $h$-weights in $y$ are $2,-2,4,0$. Interpolation in $\ad h$ isolates
$e,f,U_1+V_0,R_2$, with unit denominators at $p$. The Casimir operator
$h^2+2h+4fe$ on modules has eigenvalues $48$ on $U$ and $24$ on $V$;
it separates $U_1$ and $V_0$. Applying $e,f$ generates all of $U,V,R$,
using only the nonzero integers at most six as divisors. Their brackets
then generate $W,T,Z$ by Lemma~\ref{lem:basic}.
The modular Lie-closure computations independently obtain dimension $31$.
Equivalently, finitely many Lie words have a coefficient determinant that
is a unit in $\Zp$, and hence form a lattice basis of $B$.
\end{proof}

\section{The full finite-field automorphism group}
In this section put $k=\Fp$, and use the same letters for reductions to $k$.
The assertions about module intertwiners follow directly from
\eqref{eq:action}: the distinct weights lie between $-6$ and $6$, and the
raising/lowering coefficients are units. The indicated modules are simple,
nonisomorphic for different highest weights, and have scalar endomorphisms.

\begin{lemma}\label{lem:H1}
For $m=0,2,4,6$, one has $H^1(\mathfrak{sl}_2(k),V_m)=0$.
\end{lemma}
\begin{proof}
For a cocycle write $\alpha=c(h)$, $\beta=c(e)$, $\gamma=c(f)$.
The cocycle equations are
\[
 (h-2)\beta=e\alpha,\qquad (h+2)\gamma=f\alpha,\qquad
 \alpha=e\gamma-f\beta.
\]
Subtract a coboundary to remove the nonzero-weight components of $\alpha$.
For $m=2r>0$, the remaining component is $a v_r$. The weight-two component
of the first equation forces $ar=0$, so $a=0$. Then $\beta=bv_{r-1}$ and
$\gamma=cv_{r+1}$. The last equation gives $(r+1)(c-b)=0$, so $c=b$.
This is the coboundary of $(b/r)v_r$. All divisions are by units at $1009$.
For $m=0$, the equations immediately force $\alpha=\beta=\gamma=0$.

The independent checks give cocycle/coboundary dimension pairs
\[ (0,0),\quad (3,3),\quad (5,5),\quad (7,7) \]
for these four highest weights.
\end{proof}

Every automorphism of $\mathfrak{sl}_2(k)$ is induced by conjugation by a
matrix in $\PGL_2(k)$. For completeness, preservation of the Killing form
sends a nonzero nilpotent matrix to a nonzero nilpotent matrix. Conjugate to
fix $e$. The relations then force
\[
 h\mapsto h+a e,\qquad f\mapsto f-\frac a2h-\frac{a^2}{4}e,
\]
which is another matrix conjugation. On $V_m$, for even $m$, the compatible
$\PGL_2$ action is $\operatorname{Sym}^m(k^2)\otimes\det^{-m/2}$.
The determinant weights in \eqref{eq:transvectant} make this an action on $B$.
Denote it by $\rho$.

Let $\tau$ fix $h,e,f,U,W,T$, swap $V$ and $R$ by corresponding basis
vectors, and send $Z_0$ to $-Z_0$. The table shows that $\tau$ is an
automorphism.

\begin{proposition}\label{prop:autB}
Every automorphism of $B\otimes k$ has the form
\begin{equation}\label{eq:autB}
 \exp(\ad n)\,\rho(s)\,\tau^\epsilon,
 \qquad n\in N,\quad s\in\PGL_2(k),\quad \epsilon\in\{0,1\}.
\end{equation}
Here the exponential is a polynomial of degree at most three. In particular,
this describes the full automorphism group, including its finite components.
\end{proposition}
\begin{proof}
$N$ is the solvable radical: any solvable ideal maps to a solvable ideal in
the simple quotient $\mathfrak{sl}_2(k)$ and so lies in $N$.
After removing the quotient automorphism using $\rho$, the image of
$\mathfrak{sl}_2$ is a complement to $N$ inducing the identity on the quotient.
Its graph modulo $N'=[N,N]$ is a one-cocycle into
$N/N'\cong V_6\oplus V_4\oplus V_4$.
Lemma~\ref{lem:H1} allows conjugation by $\exp(\ad n_1)$ to remove this
graph. The remaining graph lies in the abelian ideal
$N'\cong V_6\oplus V_2\oplus V_0$ and is again a one-cocycle, removable by
$\exp(\ad n_2)$. Since $N$ has class three, the product of these inner
exponentials is another $\exp(\ad n)$; all denominators are units.
This is an explicit complement argument, not an appeal to a
positive-characteristic Levi theorem without hypotheses.

It remains to classify automorphisms $\phi$ fixing $\mathfrak{sl}_2$ pointwise.
They are equivariant. The characteristic submodule $N'$ implies
\[
 U\mapsto aU+bW,\quad W\mapsto a^2W,\quad T\mapsto cT,
 \quad Z\mapsto dZ,
\]
and the pair $(V,R)$ is mixed by an invertible two-by-two matrix $M$.
The $T$ component in $[\phi U,\phi U]$ is $2abP_5$, so $b=0$.
The bracket $[U,W]$ gives $c=a^3$. The $T$ components of the brackets on
$V\oplus R$, and the brackets with $U$, give, for $w=(1,1)$,
\begin{equation}\label{eq:block}
 M^T M=a^3 I_2,\qquad wM=a^2w.
\end{equation}
Since $M$ is invertible, also $MM^T=a^3I_2$. As $ww^T=2\ne0$,
these imply $2a^4=2a^3$, hence $a=1$.
An orthogonal matrix fixing $w$ acts by $+1$ or $-1$ on its one-dimensional
orthogonal complement. Thus $M=I$ or the swap matrix. The $Z$ bracket gives
$d=\det M$. Consequently $\phi=1$ or $\tau$.
\end{proof}

\begin{lemma}\label{lem:der}
Over $k$ and over $\Qp$, all derivations of $B$ are inner, and
$\dim\Der(B)=30$.
\end{lemma}
\begin{proof}
Let $D$ be a derivation and $\pi:B\to\mathfrak{sl}_2$ the quotient map.
The restriction $\pi D|_N$ kills $N'$ and is equivariant, since
$[D(s),n]\in N$. There is no equivariant map from
$N/N'=V_6\oplus V_4^2$ to $\mathfrak{sl}_2=V_2$; hence $D(N)\subseteq N$.
Subtract an inner derivation to make the induced derivation on the quotient
zero. The restriction to $\mathfrak{sl}_2$ is then a cocycle into $N$;
Lemma~\ref{lem:H1}, also valid over $\Qp$, removes it by another inner derivation.

For the remaining equivariant derivation write
$D(U)=aU+bW$, $D(W)=2aW$. As before $b=0$, and $D(T)=3aT$.
Writing $M$ for the action on $V\oplus R$, the linearized bracket equations are
\[
 M+M^T=3aI_2,\qquad wM=2aw.
\]
Multiplying by $w,w^T$ gives $8a=6a$, hence $a=0$.
Then $M$ is skew and annihilated by $w$, so $M=0$; also $D(Z)=0$.
Thus the remaining derivation vanishes. The centre has dimension one,
so the space of inner derivations has dimension $31-1=30$.

Independently, each implementation finds rank $931$ for the $14415$ by $961$
derivation matrix of $B$ modulo $1009$, and rank $30$ for its inner
subspace. This directly verifies the finite-field conclusion.
\end{proof}

\section{A flag with no nonidentity automorphism}
Define the flag
\begin{equation}\label{eq:flag}
 \FF_2=\langle x,y\rangle_k\subset\FF_3=\langle x,y,t\rangle_k.
\end{equation}
\begin{proposition}\label{prop:flag}
The stabilizer of \eqref{eq:flag} in the full $\Aut(B\otimes k)$ is the
identity. Its infinitesimal stabilizer in $\Der(B\otimes k)$ is zero.
\end{proposition}
\begin{proof}
Use \eqref{eq:autB}. The plane $\FF_2$ projects isomorphically to
$P=\langle h,e+f\rangle\subset\mathfrak{sl}_2$, while
$\FF_3\cap N=\langle U_0\rangle$. On $N/N'$, the inner exponential is the
identity and $\tau$ fixes $U$. Therefore $s$ preserves the highest-weight
line in $U$, so belongs to the upper Borel of $\PGL_2(k)$.
It also preserves $P$.

The Killing-orthogonal complement of $P$ is $\langle e-f\rangle$.
For a representative $\left(\begin{smallmatrix}a&b\\0&1\end{smallmatrix}\right)$,
preservation of this line forces $b=0$ and $a^2=1$. Thus $s=1$ or
$r=\operatorname{diag}(-1,1)$ projectively.

Suppose $s=r$. A flag-preserving automorphism must send $h$ to $h$ and $y$
to $-y$, since the projection of $\FF_2$ is an isomorphism.
Modulo $N'$, fixing $h$ forces $\bar n$ to have weight zero. On $U_i$, $r$
acts as $(-1)^{i+3}$, and in particular fixes $U_1$.
The $U$ component of the equation for $y$ would require
\[
 (e+f)\bar n_U=-2U_1.
\]
But $\bar n_U$ is a multiple of $U_3$, and
$(e+f)U_3=3U_2+3U_4$. This is impossible.

Now suppose $s=1$ and $\epsilon=1$. Fixing $h,y$ similarly requires
\[
 (e+f)\bar n=R_0-V_0+V_2-R_2
\]
on $V\oplus R$. The left side, applied to a weight-zero vector, has only
weights $2,-2$; the right side has nonzero weight-four and weight-zero
components. This too is impossible. The only remaining possibility is an
inner exponential fixing $x,y$. They generate $B$, so it is the identity.

For the infinitesimal assertion use Lemma~\ref{lem:der}, writing $D=\ad u$.
Its quotient component $s\in\mathfrak{sl}_2$ preserves $P$ and the highest
line in $U$, hence $s\in\langle h,e\rangle$.
The conditions $[s,h]\in P$ and $[s,e+f]\in P$ then force $s=0$.
Now $D(\FF_2)\subseteq\FF_2$ and its image in the quotient is zero.
Since $\FF_2\cap N=0$, $D(x)=D(y)=0$, and $D=0$ by generation.
\end{proof}

\section{The asymmetric lattice and its finite quotients}
\begin{definition}\label{def:A}
Use the following ordered $\Zp$-basis $b_0,\ldots,b_{30}$ of $B$:
\begin{align*}
&x,y,t,f,U_1,U_2,U_3,U_4,U_5,U_6,\\
&V_0,V_1,V_2,V_3,V_4,W_0,W_1,W_2,W_3,W_4,W_5,W_6,\\
&T_0,T_1,T_2,R_0,R_1,R_2,R_3,R_4,Z_0.
\end{align*}
It is unimodular relative to the original basis. Put
\[
 w_0=w_1=0,\quad w_2=1,\quad w_j=2\ (3\le j\le30),
 \qquad A=\bigoplus_{j=0}^{30}\Zp\,p^{w_j}b_j,
 \qquad L=p^2 A.
\]
\end{definition}
In particular $p^2B\subset A\subset B$. Although $A$ need not itself be a
Lie subring, $L$ is one: $[L,L]\subset p^4B\subset p^2A=L$.
For $i\ge1$ write $\bar L_i=L/p^iL$.
The basis $\ell_j=p^{2+w_j}b_j$ identifies its additive group with
$(\mathbb Z/p^i\mathbb Z)^{31}$. Every additive endomorphism of this
finite Lie ring is automatically $\mathbb Z/p^i\mathbb Z$-linear,
so the matrix description includes every Lie-ring automorphism.

Let $\mu$ be the bracket of $B$ and let $\nu=p^2\mu$ on the module $A$.
Under $a\mapsto p^2a$, $(A,\nu)$ is isomorphic to $L$.
Let
\[
 \EE=\End_{\Zp}(A),\qquad \EE_0=\End_{\Zp}(B),\qquad
 J=\EE\cap p\EE_0.
\]
Every endomorphism of $A$ is henceforth extended uniquely to the common
rational vector space
\[
 A\otimes_{\Zp}\Qp=B\otimes_{\Zp}\Qp.
\]
The brackets are extended by $\Qp$-bilinearity as well. Thus expressions
involving a lattice map on $B$ make sense before preservation of $B$ has
been established. For a linear map $D$, write
\[
 \delta_\mu(D)(u,v)=D\mu(u,v)-\mu(Du,v)-\mu(u,Dv),
 \qquad\delta_\nu=p^2\delta_\mu.
\]

\begin{lemma}\label{lem:lift}
Let $i\ge7$.
Every automorphism of $\bar L_i$ has a lift $Q\in1+J$.
Every solution of $\delta_\nu(D)\equiv0\pmod{p^i}$ on $A$ has a lift
$D\in J$. Here a lift means an arbitrary integral matrix representative
in $\GL(A)$ or $\EE$, respectively.
\end{lemma}
\begin{proof}
For an automorphism lift $Q\in\GL(A)$, its bracket error satisfies
\[
 E_Q(A,A)\subseteq p^{i-2}A,\qquad
 E_Q(u,v)=Q\mu(u,v)-\mu(Qu,Qv).
\]
Since $p^2B\subset A$, bilinearity gives
\begin{equation}\label{eq:error}
 E_Q(B,B)\subseteq p^{i-6}A\subseteq p^{i-6}B.
\end{equation}
This is integral when $i\ge6$. The images $Qx,Qy$ lie in $A\subset B$.
Induction on Lie words using \eqref{eq:error} now shows that $Q$ maps $B$
into $B$. By Lemma~\ref{lem:generators}, finitely many such words span $B$.
The determinant of $Q$ is a unit, independently of basis, so $Q(B)=B$.
For $i\ge7$, its reduction is an automorphism of $B/pB$.

A map preserving both $A$ and $B$ preserves
\[
 (A+pB)/pB=\FF_2,\qquad
 ((p^{-1}A\cap B)+pB)/pB=\FF_3.
\]
Proposition~\ref{prop:flag} therefore gives $Q\equiv I\pmod{p\EE_0}$.
Together with $Q-I\in\EE$, this is $Q\in1+J$.

For a derivation solution, replace $E_Q$ by $\delta_\mu(D)$.
The same word induction, now using the Leibniz error, gives $D(B)\subset B$.
Its reduction modulo $p$ is a derivation preserving the flag. The
infinitesimal part of Proposition~\ref{prop:flag} gives $D\equiv0\pmod p$
on $B$, so $D\in J$.
\end{proof}

\section{Integral logarithms count the full automorphism set}
The next step is essential: merely counting derivations or known inner
automorphisms would not determine the full automorphism group.

\begin{lemma}\label{lem:logexp}
Matrix exponential and logarithm give inverse bijections
$J\leftrightarrow1+J$. For every $i\ge3$, they respect congruence modulo
$p^i\EE$ in both directions.
\end{lemma}
\begin{proof}
The inclusions of lattices imply
\begin{equation}\label{eq:endo-bounds}
 p^2\EE_0\subseteq\EE\subseteq p^{-2}\EE_0.
\end{equation}
For $X\in J$, every $X^r$ preserves $A$, and also
$X^r\in p^r\EE_0\subseteq p^{r-2}\EE$.
In the exponential and logarithm series, denominators of terms of degree
less than $p$ are units. For terms of degree at least $p=1009$, the bounds
\[
 v_p(r!)\le r/(p-1),\qquad v_p(r)\le r/(p-1)
\]
show that all terms still lie in $\EE$ and tend to zero there.
They lie in $p\EE_0$ after removing the exponential's constant term.
Thus both series preserve the indicated sets. The usual formal inverse
identities hold in the convergent matrix algebra over $\Qp$.
In particular $1+J$ consists of invertible maps preserving $A$; its inverse
also follows from the convergent geometric series.

For precision, suppose $X,Y\in J$ and $X-Y\in p^i\EE$.
Each term in the telescoping expression for $X^r-Y^r$ lies in $p^i\EE$.
By \eqref{eq:endo-bounds} it also lies in
$p^{i+r-3}\EE_0\subseteq p^{i+r-5}\EE$.
For $r<p$, division by the relevant denominator loses no precision.
For $r\ge p$, $r-5$ exceeds both $v_p(r!)$ and $v_p(r)$.
The same conclusions therefore hold after division and summation, for
both series. Finally $p^i\EE\subset J$ when $i\ge3$, so the quotient
bijections are well defined.
\end{proof}

Put
\[
 \TT=\Hom_{\Zp}(\Lambda^2A,A),\qquad
 \TT_0=\Hom_{\Zp}(\Lambda^2B,B).
\]
The bounds $p^2B\subset A\subset B$ give
\begin{equation}\label{eq:tensor-bounds}
 p^2\TT_0\subset\TT\subset p^{-4}\TT_0,
 \quad
 p^6\End(\TT_0)\subset\End(\TT)\subset p^{-6}\End(\TT_0).
\end{equation}

\begin{lemma}\label{lem:tensor}
For $D\in J$ and $Q=\exp D$, $Q$ preserves the bracket $\nu$ modulo
$p^i\TT$ if and only if $\delta_\nu(D)\in p^i\TT$.
\end{lemma}
\begin{proof}
Let $Q$ act on bracket tensors by
$\rho_\TT(Q)\eta=Q\eta(Q^{-1}\cdot,Q^{-1}\cdot)$.
The condition $\rho_\TT(Q)\nu-\nu\in p^i\TT$ is equivalent to the
ordinary bracket-preservation condition: precompose both inputs with $Q$,
which is invertible on $A$ and preserves $p^i\TT$.
Its differential $X=\rho_{\TT,*}(D)$ is the operator
\[
 X\eta=D\eta-\eta(D\cdot,\cdot)-\eta(\cdot,D\cdot).
\]
Here $X\in\End(\TT)\cap p\End(\TT_0)$ and
$X^r\in p^{r-6}\End(\TT)$ for $r\ge6$.
Over $\Qp$, the output and two input actions commute, whence
$\rho_\TT(\exp D)=\exp X$.
Factor
\[
 \exp X-I=XU(X),\qquad
 U(X)=\sum_{r\ge0}\frac{X^r}{(r+1)!}.
\]
For $r<p-1$, the denominator is a unit. For $r\ge p-1$, the valuation
$r-6-v_p((r+1)!)$ is positive. Thus $U(X)$ preserves $\TT$.
Moreover $X^7\in p\End(\TT)$, so modulo $p$ the series for $U$ is a
polynomial in a nilpotent operator, with constant term $I$.
Consequently $U(X)$ is invertible on $\TT$. It commutes with $X$.
It follows that
\[
 (\exp X-I)\nu\in p^i\TT
 \quad\Longleftrightarrow\quad
 X\nu=\delta_\nu(D)\in p^i\TT.
\]
There is no loss of congruence precision in this equivalence.
\end{proof}

\begin{proposition}\label{prop:full-count}
For $i\ge7$, exponential and logarithm give a bijection of sets
\begin{equation}\label{eq:bijection}
 \Aut(\bar L_i)\ \longleftrightarrow\
 \ker\bigl(\delta_\nu:\EE/p^i\EE\longrightarrow\TT/p^i\TT\bigr).
\end{equation}
\end{proposition}
\begin{proof}
Every automorphism and every derivation solution lies in the respective
logarithm--exponential domain by Lemma~\ref{lem:lift}. Lemma~\ref{lem:logexp}
gives a bijection on representatives modulo $p^i\EE$, and
Lemma~\ref{lem:tensor} gives precisely the required equivalence of
nonlinear and linear bracket conditions. These statements apply to
\emph{all} automorphisms, not just to a subgroup.
\end{proof}
No assertion that the bijection is an isomorphism of additive groups is
needed or made.

\section{The exact Smith computation}
Let $c^b_{ij,k}$ be the structure constants of $B$ in the adapted basis.
In the $\ell_j=p^{2+w_j}b_j$ basis of $L$, the constants are
\begin{equation}\label{eq:scaledconstants}
 c^L_{ij,k}=p^{2+w_i+w_j-w_k}c^b_{ij,k}.
\end{equation}
All exponents are between zero and six. These constants, with $i<j$, are
also provided as ordinary integers in
\texttt{data/group\_lie\_presentation.json}.

The matrix $K$ for $\delta_\nu$ has rows $(i,j,k)$ with $i<j$, and columns
$(a,b)$ for the matrix entry $D_{ab}$. Explicitly,
\begin{equation}\label{eq:matrix}
 K_{(i,j,k),(a,b)}
 =\mathbf1_{k=a}c^L_{ij,b}
  -\mathbf1_{i=b}c^L_{aj,k}
  -\mathbf1_{j=b}c^L_{ia,k}.
\end{equation}
Thus $K$ is an integer matrix of shape $14415\times961$.

\begin{proposition}[Exact computational certificate]\label{prop:smith}
The rank of $K$ over $\Qp$ is $931$. Its nonzero $p$-primary Smith
valuations are
\begin{center}
\begin{tabular}{rrrrrr}
\toprule Valuation & 0 & 1 & 2 & 3 & 4 \\
Multiplicity & 87 & 55 & 758 & 6 & 25 \\
\bottomrule
\end{tabular}
\end{center}
Their sum is $1689$; the nullity is $30$.
\end{proposition}
\begin{proof}[Certificate and verification]
The Python implementation constructs the integer brackets from
\eqref{eq:falling}, changes basis, builds \eqref{eq:matrix}, and performs
local Smith elimination modulo $p^{12}$. Each pivot has minimum $p$-valuation
in the remaining block. Only unit inverses are taken; divisions by a
pivot power of $p$ are checked to be exact. The resulting plan consists
of $931$ original row indices, column indices and valuations.

A separately written C++ implementation independently reconstructs the
brackets, the basis transformation and the entire matrix from the formulas.
It does not read the Python coefficient tables. It verifies the proposed
plan modulo $p^6$, checking divisibility of the active block at each new
valuation level, pivot normalization, exact elimination, and a zero residual
block. It finds the same five multiplicities. The Python implementation
also replays the plan at precision six.

Every reported valuation is less than six, so it is determined exactly by
this finite-ring calculation. Integer Jacobi places the inner derivations
of $B$ in its derivation kernel. Their verified rank $30$ modulo $p$ supplies
thirty that remain linearly independent over $\Qp$. An invertible change of basis
transports them by conjugation, and multiplying the bracket by $p^2$ does
not change its rational derivation kernel, since
$\delta_{p^2\mu}=p^2\delta_\mu$. They therefore yield $30$ independent
vectors in the rational kernel for $L$, proving rank at most $961-30=931$.
The $931$ nonzero pivots give rank at least $931$. There can therefore be
no additional nonzero invariants hidden at higher precision. All matrix
operations are exact elementary operations over the indicated local ring.
The sum is $55+2(758)+3(6)+4(25)=1689$.
\end{proof}

\begin{corollary}\label{cor:autLie}
For every $i\ge7$,
\begin{equation}\label{eq:autLie}
 |\Aut(\bar L_i)|=p^{30i+1689}.
\end{equation}
\end{corollary}
\begin{proof}
In Smith coordinates, each of the $30$ zero diagonal entries contributes
$p^i$ kernel elements; each nonzero entry $p^v$ contributes $p^v$ because
$i>4\ge v$. Combine Proposition~\ref{prop:smith} with the full-set
bijection \eqref{eq:bijection}.
\end{proof}

\section{The finite group and the equality}
Now set $i=1689$ and define
\begin{equation}\label{eq:G}
 G=\BCH(\bar L_{1689}).
\end{equation}
This is an explicit construction, not a library identifier: its elements
are the $31$-tuples modulo $1009^{1689}$ in the $\ell$ basis, with the Lie
bracket \eqref{eq:scaledconstants}. Its group product is the BCH Lie
polynomial truncated after bracket length $846$.

\begin{lemma}\label{lem:class}
The Lie ring $\bar L_{1689}$ has order $p^{52359}$ and nilpotency class
at most $846$.
\end{lemma}
\begin{proof}
$L$ is free of rank $31$ over $\Zp$, so its quotient has precisely
$p^{31\cdot1689}=p^{52359}$ elements, with unique additive coordinates.
Since $L\subset p^2B$ and $B$ is a Lie lattice,
$\gamma_m(L)\subset p^{2m}B$ for every $m$.
On the other hand $p^{1689}L=p^{1691}A$ contains $p^{1693}B$.
Consequently
\[
 \gamma_{847}(L)\subset p^{1694}B\subset p^{1693}B
 \subset p^{1689}L,
\]
which proves the class bound.
\end{proof}

For clarity, BCH here can be defined without any software convention.
Take the terms of degree at most $846$ in
$\log(\exp X\exp Y)$ in the completed free associative algebra over
$\mathbb Q$ and express each homogeneous term as a Lie polynomial.
For example, the right-normed Dynkin operator divided by the degree
provides such an expression. Evaluate these Lie polynomials using the
displayed bracket. The denominators have only prime factors at most $846$,
so all coefficients have unique interpretations modulo $1009^{1689}$.
This is a finite, canonical rule for multiplication; it need not be expanded
into a very large word list to specify the group.

The Lazard correspondence applies to finite $p$-Lie rings of class strictly
less than $p$ \cite{Lazard,CDV}. It preserves the underlying set, group order,
and the full automorphism group; see in particular the discussion on
pages 433--434 of \cite{CDV}. Thus the BCH product is associative, has
identity zero and inverse $-u$, and every group automorphism corresponds
to a Lie-ring automorphism. No exponent-$p$ hypothesis is required.
The actual class bound is $846<1009$, so these are the applicable
hypotheses, rather than a general assertion beyond the Lazard range.

\begin{proof}[Proof of Theorem~\ref{thm:main}]
Lemma~\ref{lem:class} proves the exact order and validates the BCH construction.
Corollary~\ref{cor:autLie} and the full-automorphism preservation give
\[
 |\Aut(G)|=|\Aut(\bar L_{1689})|
 =1009^{30\cdot1689+1689}
 =1009^{52359}=|G|.
\]
The prime is odd and the group is nontrivial. This answers the stated
existence question affirmatively.
\end{proof}

\section{Reproducibility and verification boundaries}\label{sec:reproducibility}
\noindent\textbf{Source and certificates.}
Version 1.1.0 has archive DOI
\href{https://doi.org/10.5281/zenodo.22770864}{\nolinkurl{10.5281/zenodo.22770864}}.
The matching manuscript, explicit group specification, exact certificates,
verification software, and audit records are available in the fixed
publication snapshot:
\begin{center}
\href{https://github.com/AlecKriebel/Math/tree/kourovka-16-63-v1.1.0/kourovka_16_63}{\texttt{kourovka-16-63-v1.1.0}}.
\end{center}
The \href{https://aleckriebel.github.io/Math/papers/kourovka-16-63/}{project landing page}
provides the PDF and companion archive. The archive includes a
\texttt{SHA256SUMS} integrity manifest. The fixed version tag identifies this
edition independently of future changes to the development branch.

\subsection{Fast independent verification}
From the root of the bundle, run
\begin{verbatim}
python3 scripts/resume.py
\end{verbatim}
This compiles and runs the independent C++ verifier and also tests that a
corrupted pivot certificate is rejected. Alternatively, with a GCC or
Clang compiler supporting 128-bit integers, run
\begin{verbatim}
mkdir -p build
c++ -std=c++17 -O2 -o build/independent_verify \
    scripts/independent_verify.cpp
./build/independent_verify data/smith_pivot_plan.txt
\end{verbatim}
The C++ verification uses only the standard library and the supplied
pivot plan. It reconstructs all mathematical input coefficients itself.
No GAP, SageMath, Magma, remote computation, or randomized group
identification is needed for this check.

\subsection{Regeneration}
The independent Python construction requires NumPy and SymPy. The exact
commands are
\begin{verbatim}
python3 scripts/lie31.py
python3 scripts/exact_linear.py
python3 scripts/smith_local.py --precision 12
python3 scripts/smith_local.py --verify-plan --precision 6
python3 scripts/export_presentation.py
python3 scripts/sanity_groups.py
\end{verbatim}
The bundle includes executed logs, the coefficient tables, both independent
implementations, the Smith pivot plan, and the unscaled rank certificate.
The original supplied run used Python 3.13.5, NumPy 2.3.5, SymPy 1.14.0,
and g++ 14.2.0 (Debian 14.2.0-19). Independent full regeneration for this
edition used Python 3.12.14, NumPy 2.3.5, SymPy 1.14.0, and Apple Clang
21.0.0 on macOS 26.6.2 (arm64). The only verifier repair was an
indentation-only portability correction for Clang. The strict build used
\texttt{-std=c++17 -O2 -Wall -Wextra -Werror}.
The C++ local modulus was the exact integer
$1009^6=1055229678769825441$; products use 128-bit arithmetic.
The measured times in each log describe its recorded environment only.
Fresh records are in \texttt{audit/reproduction/}. A third implementation
using arbitrary-precision Python integers independently reconstructs the
brackets, checks the exported presentation and a $931$-square rank minor,
and verifies the Smith data at precision five. A separate implementation
checks the finite-field generation, involution, cohomology, and flag ranks.
The C++ verifier also passes an address- and undefined-behavior-sanitized
build. The independent audit reports record the scope of each check.

\subsection{What the computations do and do not certify}
The computations verify integer Jacobi, generation, cohomology dimensions,
derivation and flag ranks, and all local Smith invariants. The full
finite-field stabilizer, lifting, integral logarithm--exponential, and Lazard
steps are proved above; the numerical output alone does not certify them.

Neither the group's Cayley table nor its full automorphism set is enumerated.
The canonical BCH rule through bracket length $846$ and the explicit Lie table
specify multiplication without expanding every monomial. Independent
implementations and scoped AI reviews provide reproducible evidence; they
are not external human peer review or proof-assistant formalization.

\clearpage
\appendix
\section{Selected inspected output}
The final independent C++ run includes
\begin{verbatim}
original_integer_jacobi_checks=4495
adapted_integer_jacobi_checks=4495
unscaled_derivation_rank_mod_p=931 derivation_dimension=30
smith_valuation=0 multiplicity=87
smith_valuation=1 multiplicity=55
smith_valuation=2 multiplicity=758
smith_valuation=3 multiplicity=6
smith_valuation=4 multiplicity=25
sum_smith_valuations=1689
candidate_group_order_exponent=52359
predicted_aut_order_exponent=52359
ALL_CHECKS_PASSED
\end{verbatim}
Some long log lines are split here for readability. The names
\texttt{candidate\_group\_order\_exponent} and
\texttt{predicted\_aut\_order\_exponent} are retained from the computational
checkpoint; the proof establishes the corresponding group statements.
A separate negative test changes the first proposed pivot valuation from
zero to five. The verifier exits nonzero with a pivot-valuation mismatch.

\begin{thebibliography}{9}
\bibitem{KN}
E. I. Khukhro and V. D. Mazurov (eds.),
\emph{The Kourovka Notebook: Unsolved Problems in Group Theory},
21st edition, update of 1 September 2026. Problem 16.63, p. 102;
Archive 12.77, p. 255.
\href{https://kourovkanotebookorg.wordpress.com/wp-content/uploads/2026/09/21tkt.pdf}{September 2026 edition (PDF)}.

\bibitem{GSJ}
J. Gonz\'alez-S\'anchez and A. Jaikin-Zapirain,
\emph{Finite $p$-groups with small automorphism group},
Forum of Mathematics, Sigma 3 (2015), e7.
\url{https://arxiv.org/abs/1406.5772}.

\bibitem{SA}
M. Shabani-Attar,
\emph{On equality of order of a finite $p$-group and order of its automorphism group},
Bull. Malaysian Math. Sci. Soc. 38 (2015), 461--466.
\url{https://doi.org/10.1007/s40840-014-0030-z}.
Also available as the
\href{https://math.usm.my/bulletin/pdf/acceptedpapers/2012-10-027-R1.pdf}{author-accepted manuscript (PDF)}.

\bibitem{OR}
B. Omirov and J. Ruan,
\emph{On Lie Algebras with Only Inner Derivations},
arXiv:2605.04602v1 (6 May 2026), Theorem 4.8.
\url{https://arxiv.org/html/2605.04602v1}.

\bibitem{Lazard}
M. Lazard,
\emph{Sur les groupes nilpotents et les anneaux de Lie},
Ann. Sci. \'Ecole Norm. Sup. (3) 71 (1954), 101--190.
\href{https://www.numdam.org/item/?id=ASENS_1954_3_71_2_101_0}{NUMDAM archive}.

\bibitem{CDV}
S. Cical\`o, W. A. de Graaf and M. Vaughan-Lee,
\emph{An effective version of the Lazard correspondence},
J. Algebra 352 (2012), 430--450, especially pp. 430, 433--434.
\url{https://doi.org/10.1016/j.jalgebra.2011.11.031}.
\href{https://iris.unitn.it/retrieve/e3835192-3015-72ef-e053-3705fe0ad821/lazard.pdf}{author-hosted manuscript (PDF)}.
\end{thebibliography}
\end{document}
